\subsection{Posibles sesgos y fuga de información}

La documentación presenta los datos como un ejemplo de una empresa ficticia~\cite{ibm}. No se conoce 
cómo se seleccionaron los clientes, por lo que no puede afirmarse que la muestra represente al mercado 
peruano. Además, algunos perfiles tienen menor presencia: \var{SeniorCitizen=1} reúne 1142 clientes, 
frente a 5901 en la otra categoría. Por estas diferencias de tamaño y de abandono, será útil revisar 
los errores del modelo por grupo. Sin embargo, estas diferencias por sí solas no demuestran discriminación.

También existe una limitación temporal. Los clientes de mayor antigüedad son quienes permanecieron 
el tiempo suficiente para alcanzar ese grupo. Compararlos con clientes recientes no permite saber qué 
pasaría con una misma persona si permaneciera más tiempo. Además, no se cuenta con las fechas en que 
se registró cada variable, por lo que falta comprobar si esa información estaba disponible antes del 
abandono.

La fuga de información ocurre cuando el entrenamiento utiliza información que no estaría disponible al 
predecir~\cite{leakage}. Se excluirán de las entradas \var{Churn} y su copia numérica \var{abandono}, 
pues ambas representan el objetivo. \var{customerID} también se excluirá, ya que identifica a cada 
cliente y no describe sus características. Las correlaciones entre cargos y antigüedad pueden indicar 
información redundante, pero no son por sí solas una fuga de información.


\subsection{Síntesis del EDA y siguientes pasos}

El EDA identifica diferencias por contrato, antigüedad, internet, método de pago y ciertos complementos. 
Las comparaciones conjuntas muestran que algunas diferencias permanecen al separar por contrato o 
antigüedad. No establecen causas ni determinan todavía qué variables aportarán más a un modelo.

Para P2 se propone definir el momento de predicción y separar los datos de entrenamiento y prueba 
manteniendo la proporción de clases. La imputación y la codificación se ajustarán únicamente con 
los datos de entrenamiento. Se conservará inicialmente la antigüedad original, omitiendo el 
grupo creado solo para visualizar. Se comparará un modelo que siempre prediga la clase mayoritaria 
con al menos dos modelos supervisados, por ejemplo regresión logística y un árbol de decisión. 
Sus parámetros se elegirán con los datos de validación, sin utilizar los de prueba. 
Se reportarán matriz de confusión, precisión y recall, junto con una métrica para evaluar 
qué tan bien el modelo asigna mayor riesgo a quienes abandonan. 
